\section{Introduction}

In recent years a number of studies concerning computational complexity in video games has been conducted. The studies have been on all kinds of genres, from platformers to puzzle games. So far these studies seem to have ignored the popular puzzle game Lumines. In this report a simplified model of Lumines is constructed. A decision problem is then formulated in the model. The decision problem is proven to be in NP. A reduction from the NP-complete \textit{Subset sum} to the game model is then presented, showing that the simplified version of Lumines is NP-complete. 

Lumines is a game of interest because to the authors' knowledge it is not well researched from the perspective of computational complexity. Several similar games has been more thoroughly researched in recent years and comparing Lumines to these findings might yield interesting results. New findings on the computational complexity of Lumines may aid in the development of video game artificial intelligence. Our understanding of computational complexity in generic 2-dimensional grid games may also benefit from this research. Combining results from this report with results from previous similar research might make it possible to find a correlation between certain game characteristics and classes of computational complexity.

\subsection{Technical Overview}

An environment modeling Lumines is formulated and the computational complexity of this environment is shown to be NP-complete. In the domain of computational complexity, Lumines has not been researched priorly to this report. The proof in this report uses a reduction from the \textit{subset sum} problem. \textit{Subset sum} has not been used extensively in previous studies on computational complexity in video games.

\subsection{The game}
The video game Lumines was released in December 2004 in Japan for the PSP platform. The game was a major success for the developer Q Entertainment, and has since its original release been followed with a couple of sequels \cite{wiki:lumines}.
\begin{figure}[H]
    \centering
    \resizebox{0.7\textwidth}{!}{
    \begin{tikzpicture}
        \lummonowhite{0}{0}
        \lummonoblack{3}{0}
        \lumlwhite{6}{0}
        \lumlblack{9}{0}
        \lumh{12}{0}
        \lumx{15}{0}
    \end{tikzpicture}
    }
    \caption{The six pieces of lumines: M\{W,B\}, L\{W,B\}, H and X}
    \label{fig:pieces}
\end{figure}

Lumines is a two-dimensional grid game, in which a sequence of $2 \times 2$ dichromatic \textit{blocks} fall The player rotates and shifts each block before its bottom \text{cells} reaches other cells, \textit{terrain}, in the gameboard. When the block reach existing terrain, it is no longer possible to shift or rotate it. If both bottom cells of the block is supported by existing terrain, the block stays put on that terrain. Otherwise, the half of the block that is not supported by terrain will fall down until it is. The player is then presented with the next block.

\begin{figure}[H]
    \centering
    \resizebox{0.6\textwidth}{!}{
    \begin{subfigure}[b]{0.45\linewidth}
        \resizebox{\linewidth}{!}{
        \begin{tikzpicture}
            \lumx{0}{4}
            \draw[thick, ->] (1,3.5) -- (1, 2.5);
            \lumx{0}{0}
            \draw [draw=none] rectangle (3, 7);
        \end{tikzpicture}
        \hspace{0.25cm}
        \begin{tikzpicture}
            \lumx{0}{2}
            \lumx{0}{0}
            \draw [draw=none] rectangle (3, 7);
        \end{tikzpicture}
        }
        \caption{}
    \end{subfigure}
    \hspace{0.05\linewidth}
    \begin{subfigure}[b]{0.45\linewidth}
        \resizebox{\linewidth}{!}{
        \begin{tikzpicture}
            \lumx{1}{4}
            \draw[thick, ->] (2,3.5) -- (2, 2.5);
            \lumx{0}{0}
            \draw [draw=none] rectangle (3, 7);
        \end{tikzpicture}
        \hspace{0.25cm}
        \begin{tikzpicture}
            \lumx{0}{0}
            \cellb{1}{2}
            \cellw{1}{3}
            \cellw{2}{0}
            \cellb{2}{1}
            \draw [draw=none] rectangle (3, 7);
        \end{tikzpicture}
        }
        \caption{}
    \end{subfigure}
    }
    \caption{Block behavior with and without supporting terrain}
\end{figure}

If at any point the player forms a monochromatic $2 \times 2$ square, all cells belonging to that square is \textit{marked}. Note that these squares may overlap, for example all cells within a monochromatic $2 \times 3$ rectangle is marked, since it consists of overlapping $2 \times 2$ monochromatic squares. At a regular pace a vertical \textit{sweep-line} scans the gameboard and removes all marked cells.

\begin{figure}[H]
    \centering
    \resizebox{0.6\textwidth}{!}{
        \begin{subfigure}[b]{.3\linewidth}
        \fbox{
        \begin{tikzpicture}
            \cellw{0}{0}
            \cellw{1}{0}
            \cellb{2}{0}
            \cellw{0}{1}
            \cellw{2}{1}
            \cellw{2}{2}
            \cellb{0}{5}
            \cellw{1}{5}
            \cellb{0}{6}
            \cellw{1}{6}
            \draw [draw=none] (0,0) rectangle (3, 8);
        \end{tikzpicture}
        }
        \end{subfigure}
        \begin{subfigure}[b]{.3\linewidth}
        \fbox{
        \begin{tikzpicture}
            \cellw{0}{0}
            \cellw{1}{0}
            \cellb{2}{0}
            \cellw{0}{1}
            \cellw{2}{1}
            \cellw{2}{2}
            \cellb{0}{2}
            \cellw{1}{1}
            \cellb{0}{3}
            \cellw{1}{2}
            \draw [draw=none] (0,0) rectangle (3, 8);
        \end{tikzpicture}
        }
        \end{subfigure}
        \begin{subfigure}[b]{.3\linewidth}
        \fbox{
        \begin{tikzpicture}
            \cellb{0}{2}
            \cellb{0}{3}
            \cellb{2}{0}
            \draw (0,0) [very thick, draw=black, fill=red] rectangle (2,2);
            \draw (1,1) [very thick, draw=black, fill=red] rectangle (3,3);
            \draw (1,1) [very thick, draw=black, fill=red] rectangle (2,2);
            \draw [draw=none] (0,0) rectangle (3, 8);
        \end{tikzpicture}
        }
        \end{subfigure}
    }
    \caption{Marking of $2 \times 2$ monochromatic squares \cite{lumines})}
\end{figure}

The version of Lumines found in the actual video games is played on a $10 \times 16$-gameboard. The goal of the standard game mode in Lumines is to survive, to keep placing blocks inside the gameboard, for as long as possible and clear as much terrain as possible in the process. The blocks given to the player is generated randomly, and 3-block lookahead is provided. The game starts on a partially filled gameboard, which varies in size depending on the difficulty level \cite{lumines}.

\subsection{Online and offline Lumines}
In Lumines all playable blocks are generated probabilistically. In this report this feature will be refered to as \textit{online}. In contrast, a game of Lumines where all playable blocks are known beforehand will be referred to as \textit{offline}.

\subsection{Problem statement}
This report considers the generalized version of the game, which is played on an $n \times m$-gameboard. The authors seek to formalize and simplify Lumines, and explore this new system's computational complexity characteristics. The maximization of cleared cells is the main goal this paper examines and is translated into a decision problem ``Can k cells be cleared given a particular Lumines game?''. 

While the sweep-line is a core characteristic of Lumines, all marked terrain will be considered instantly cleared. Furthermore only the offline version of the game with no allowed block rotations during gameplay will be considered. These are the three main simplifications  assumed in this report, although a more thorough discussion is provided in \nameref{method}. The impact of these simplifications on the validity of the results is considered in \nameref{discussion}.

The main problem this report intend to answer is:
\begin{itemize}
        \item Is (standard mode) Lumines NP-complete?
\end{itemize}
